\documentclass[border=3pt]{standalone}
% 公共导言区 —— 所有 TikZ 图 \input 本文件后使用
\usepackage[UTF8, fontset=mac]{ctex}
\usepackage{tikz}
\usetikzlibrary{arrows.meta, positioning, calc, shapes.geometric, decorations.pathmorphing, backgrounds, fit}
\definecolor{zhusha}{HTML}{9B2D20}
\definecolor{mohei}{HTML}{1A1A1A}
\definecolor{shenhui}{HTML}{555555}
\definecolor{qianhui}{HTML}{F5F0EC}
\definecolor{lan}{HTML}{1F4E79}
\definecolor{qing}{HTML}{2E7D6E}
\tikzset{
  block/.style={draw=shenhui, fill=white, thick, rounded corners=1.5pt,
                font=\small, align=center, inner sep=4pt, minimum height=7mm},
  core/.style={block, fill=lan!12, draw=lan, font=\small\bfseries},
  periph/.style={block, fill=qing!10, draw=qing},
  power/.style={block, fill=zhusha!8, draw=zhusha},
  note/.style={font=\footnotesize, text=shenhui, align=center},
  lab/.style={font=\footnotesize, text=mohei},
  sig/.style={-{Stealth[length=2.2mm]}, thick, color=mohei},
  fbsig/.style={-{Stealth[length=2.2mm]}, thick, color=zhusha, dashed},
  vec/.style={-{Stealth[length=3mm]}, very thick, color=zhusha},
}

\begin{document}
\begin{tikzpicture}[x=1mm, y=1mm, font=\small, >=Stealth]

% ===== 六边形顶点（基本矢量）=====
\def\R{34mm}
\foreach \k in {1,...,6} {
  \pgfmathsetmacro{\a}{(\k-1)*60+30}
  \coordinate (V\k) at (\a:\R);
}
% 六边形的 6 条边
\draw[very thick, lan] (V1) -- (V2) -- (V3) -- (V4) -- (V5) -- (V6) -- cycle;

% 基本矢量箭头 + 标签（开关状态）
\foreach \k/\st/\labpos in {1/100/{right}, 2/110/{above right}, 3/010/{above left}, 4/011/{left}, 5/001/{below left}, 6/101/{below right}} {
  \draw[vec, color=lan] (0,0) -- (V\k);
  \node[\labpos, lab, text=lan] at (V\k) {$\vec{V}_{\k}$\\{\scriptsize (\st)}};
}

% ===== 扇区标注 =====
\foreach \s/\a in {I/30, II/90, III/150, IV/210, V/270, VI/330} {
  \node[lab, text=shenhui] at (\a:11mm) {\s};
}

% ===== 零矢量 =====
\node[circle, fill=zhusha, inner sep=1.2pt] at (0,0) {};
\node[note, anchor=north east, text=zhusha] at (-1mm, 1mm) {$\vec{V}_0(000)$\\$\vec{V}_7(111)$};

% ===== 参考矢量（扇区 I）=====
\coordinate (ur) at (50:{0.86*\R});
\draw[vec, very thick] (0,0) -- (ur);
\node[above right=1mm, lab, text=zhusha] at (ur) {$\vec{U}_{ref}$（角度 $\theta$）};

% 平行四边形分解：T1 沿 V1，T2 沿 V2
\coordinate (p1) at ($(V1)!0.5!(ur)$);
\draw[dashed, thick, zhusha] (ur) -- ++($(-30:0.0)$);
\path (ur) -- ++(210:{0.42*\R}) coordinate (h1);
\draw[dashed, thick, zhusha] (ur) -- (h1) -- (V1);
\node[note, text=zhusha] at ($(ur)!0.5!(h1)$) [above left=0.3mm] {$T_2/T$};
\node[note, text=zhusha] at ($(V1)!0.5!(h1)$) [below=0.4mm] {$T_1/T$};
\node[note, text=zhusha] at ($(0,0)!0.5!(V1)$) [below right=0.3mm] {$T_0/T$（零矢量）};

% ===== 线性区圆 =====
\draw[thin, qing, dashed] (0,0) circle (29.4mm);
\node[note, text=qing, anchor=south west] at (150:30mm) {线性区边界 $|\vec{U}| \le V_{dc}/\sqrt{3}$};

% ===== 幅值标注 =====
\node[note, anchor=south] at (90:{\R+3mm}) {基本矢量幅值 $2V_{dc}/3$，互差 60°};

\end{tikzpicture}
\end{document}
